% !PRINTED 160
\weberpage{160}
%<<CONT
die höchste Potenz von $y$ aus der Endgleichung wegfällt. Man
würde dann die Uebereinstimmung der Zahl mit (5) nur durch
die willkürliche Hinzufügung unendlicher Wurzeln retten können.
Auch mehrfache Wurzeln geben zu Bedenken Anlass. Man begegnet diesen Uebelständen theilweise dadurch, dass man die
% [corrected 2026-08-07, 6x zoom strip 2/6] "homogene" alone is letterspaced -- the term at the
% point of definition. "Form der Gleichungen zu Grunde legt, eine Form," is set normally, so the
% run stops at the one word. Claimed by the translation pass and, independently and on a
% different model, by the sperrung pass; both correct. The same page's "Uebelstände" query is
% REFUTED -- the fragment already reads "Uebelständen", as the page does.
\gesperrt{homogene} Form der Gleichungen zu Grunde legt, eine Form,
in der sich das Eliminationsproblem besonders in der Geometrie
einstellt.

Wenn nämlich die Function $f(x)$ aus einer homogenen
Function $n^{\mathrm{ter}}$ Ordnung der drei Variablen $x, y, z$ (einer ternären
Form) durch Ordnen nach Potenzen von $x$ hervorgegangen ist,
so sind die Co\"efficienten $a_0, a_1, a_2 \ldots a_n$ homogene Formen der
beiden Variablen $y, z$ von dem Grade, den der Index angiebt,
also $a_0$ eine Constante, $a_1$ eine lineare, $a_2$ eine quadratische
Form u.\,s.\,f., und Entsprechendes gilt von den Co\"efficienten
$b_0, b_1 \ldots b_m$ der Function $\varphi(x)$.

Bedeuten $x, y, z$ Dreieckscoordinaten in der Ebene, so sind
$f(x) = 0$, $\varphi(x) = 0$ die Gleichungen einer Curve $n^{\mathrm{ter}}$ und
$m^{\mathrm{ter}}$ Ordnung, und die gemeinsamen Wurzeln dieser beiden
Gleichungen (die Verhältnisse $x:y:z$) bedeuten die Schnittpunkte
der beiden Curven. Die Endgleichung, die man durch Elimination von $x$ erhält, ist eine homogene Gleichung, deren Grad
sich aus (15), \S.\,48 gleich $nm$ ergiebt.

Eine Verminderung des Grades kann hier nicht stattfinden;
immer ist die Endgleichung eine homogene Gleichung des $nm^{\mathrm{ten}}$
Grades für die beiden Unbekannten $y, z$. Die einzige Ausnahme,
auf die hier zu achten ist, ist die, dass die Resultante identisch
(für alle $y, z$) verschwindet, dass also unendlich viele gemeinsame Wurzeln vorhanden sind. Dieser Fall hat in der Geometrie
die Bedeutung, dass die beiden Curven einen Curventheil und
nicht bloss einzelne Punkte gemein haben.

Wenn die Gleichungen $f = 0$, $\varphi = 0$ nur eine endliche
Anzahl gemeinsamer Wurzeln haben, so können wir in den
Ausdrücken

\[
y_1 = y + \beta x, \qquad z_1 = z + \gamma x
\]

die Constanten $\beta, \gamma$ so einrichten, dass, wenn einem Werth des
Verhältnisses $y : z$ mehrere verschiedene der Gleichungen (1)
genügende Werthe von $x$ entsprechen, die zugehörigen Werthe
von $y_1 : z_1$ von einander verschieden ausfallen.
